% Appendix 15: Scenario Omicron - Kinship \\\& Household Interfaces
\section{Appendix 15: Scenario Omicron - Kinship \\\& Household Interfaces}

\subsection{The Legacy Paradigm \\\& Its Failures}
Legacy models treat households as collections of disparate individual accounts or clumsy "family plans" tied to a single centralized billing entity. Sharing access to smart home devices, finances, or scheduling involves fractured permissions, high friction, and relying on centralized cloud servers that break during outages.

\subsection{The Incremental Automation Pathway}
\textbf{Phase 1: Manual Sharing} - Users explicitly grant access to local devices or communal resources via QR codes or local Wi-Fi handshakes.
\textbf{Phase 2: Ambient Assistance} - UIs suggest resource sharing based on proximity and historical usage, presenting low-friction approval prompts.
\textbf{Phase 3: Ephemeral Trust} - Full zero-click interaction where the Oasis Engine infers kinship trust boundaries through continuous local mesh interactions, seamlessly allocating shared bandwidth, power, and computational resources.

\subsection{The Human Boundary (Limits of Automation)}
The system must never automate the addition or removal of a member from the core kinship circle. These events (e.g., adding a new dependent, or severing a compromised device from a household) require a high-friction, multi-party confirmation using physical presence or biometric verification.

\subsection{Interface Mechanics (The "How")}
Household state is maintained via a decentralized CRDT lattice, entirely bypassing the cloud. Local devices act as a resilient mesh. The UI relies heavily on ambient E-ink displays within the household for shared state (e.g., microgrid battery levels, communal task completion), minimizing the need for individuals to pull out their personal screens. Trust boundaries are visually represented through spatial AR overlays showing which resources are currently pooled versus private.

\subsection{Semantic NLP Translation}
When a user expresses a messy intent like "Let my sister use the solar battery for her car," the local LLM parses this natural language against the household's CRDT state. It identifies "sister" and "solar battery," resolving these to specific cryptographic identities and resources. The LLM then formulates a strict Layer 5 \texttt{ResourceDelegationIntent} payload, defining the exact power limits and duration, which is presented to the user for final cryptographic signature.
